\section{A society of language-model agents}
\label{sec:llm}

The phase diagram of \S\ref{sec:results} is a claim about agents that learn from
one another under a trust-weighted rule. Perceptrons instantiate that rule
exactly; language-model agents instantiate something we cannot write down. If the
phase structure survives the substitution, it is a property of the interaction
rather than of the parameterisation.

$N$ language-model agents share an agenda of $P$ issues and interact in rounds
mirroring \S\ref{sec:model}: an issue, an emitter and a receiver are drawn; the
emitter states its position with a brief justification; the receiver, seeing the
statement and the emitter's group label, updates both its own position on that
issue and its disposition towards that emitter. Nothing instructs an agent to
agree with, trust, or favour anybody. The two axes of the phase diagram become two
knobs on the instructions: $\fd$ is the fraction of agents whose prompt carries a
discriminating clause, assigned independently of group, and $d$ is the strength of
that clause --- a graded instruction to treat the other group's statements as less
credible than their content warrants and one's own group's as more credible, at
several levels including $d=0$ (clause absent) and $d<0$ (the clause reversed).
Group labels name arbitrary partitions with no content relating to any issue, so
that class-aligned structure has to be produced by the interaction rather than
retrieved from what the model already associates with a real social category.
Appendix~\ref{app:llm-prompts} gives the protocol in full.

Measurement needs only the two quantities of \S\ref{sec:order-params}, both
elicited outside the interaction loop so that measuring does not perturb the
dynamics. Each agent rates all $P$ issues on $[-1,1]$, giving an opinion vector
$v_I$ whose normalised inner products are the overlaps $\rho_{IJ}$; and each agent
reports, for every other agent, how far it trusts that agent's statements, mapped
to $\eta_{J|I}\in[-1,1]$ and deliberately not symmetrised. The three correlations
\eqref{eq:correlations} and the two balance aggregates \eqref{eq:balance} then
follow unmodified, which is the point of having defined them without reference to
weights or to $K$.

\subsection{Results}
\label{sec:llm-results}

\llmtodo{To be written from the run. The comparisons it must make, in order of
importance: (i) the sign and magnitude of $\Rmuc$ against $d$ at fixed $\fd$,
versus the sharp transition at $d\approx0$ of Figure~\ref{fig:correlations};
(ii) whether the two discriminatory regimes are distinguishable, i.e. whether
$\Rcw$ falls while $\Rmuc$ stays high as $d$ grows; (iii) whether $d<0$ gives the
frustrated signature $\BA<0$ with $\BI\approx0$ rather than simple reverse
discrimination; (iv) the agenda-complexity reversal.}

% ---------------------------------------------------------------------------
% RESERVED SLOT for the language-agent figure.  Commented out while the
% panels hold no data: a figure with nothing measured in it does not earn
% space in a 9-page limit.  Uncomment when the study runs, and take the
% room back from the reserve list in NOTES.md.
% ---------------------------------------------------------------------------
% \begin{figure}[t]
%   \centering
%   \begin{subfigure}{0.55\linewidth}
%     \includegraphics[width=\linewidth]{figures/placeholder_llm_correlations}
%     \caption{Order parameters versus instructed intolerance.}
%   \end{subfigure}\hfill
%   \begin{subfigure}{0.33\linewidth}
%     \includegraphics[width=\linewidth]{figures/placeholder_llm_balance}
%     \caption{Balance trajectories.}
%   \end{subfigure}
%   \caption{\textbf{The same measurements in a society of language-model agents.}
%   \llmtodo{Placeholder. (a) will show $\Rmuc$, $\Rcw$ and $\Rwmu$ against the
%   instructed field $d$ at fixed $\fd$, overlaid on the corresponding cut through
%   the simulated phase diagram. (b) will show $(\BI,\BA)$ trajectories for a simple
%   and a complex agenda, against Figure~\ref{fig:trajectories}.}}
%   \label{fig:llm}
% \end{figure}


A match would establish the phase structure as a property of trust-weighted social
learning as such --- and, the transition being sharp, that a small change in how
agents are told to treat each other can move a deployed population between regimes.
A mismatch is equally informative, and we commit in advance to reporting which of
the four comparisons fails. The likely failure modes are statements about the
substrate rather than the model: stated trust compressed towards the positive end,
instruction drift making $d$ non-stationary, or an agenda too small to reach
$\alpha>1$.
